%! $n$th Roots \& Rational Exponents — Unit 6, Lesson 6.1 — Group Activity (Answer Key)
\documentclass[10pt]{article}
\usepackage{algebra2-article}
\usepackage{algebra2-key}   % pulls in -boxes; do NOT also load -boxes
\newcommand{\TallMath}[1]{$\displaystyle #1\rule[-1.4em]{0pt}{3.2em}$}
% In-formula answer slot (\ans is text-mode and must never appear inside $...$).
\newcommand{\mans}[1]{{\color{keyred}\mathbf{#1}}}

\begin{document}

\pageheader{Unit 6, Lesson 6.1}{Group Activity --- Answer Key}
\namepartnerperiod

\vspace{0.08in}

\begin{headlinebox}{forestbg}
\small\textbf{One Expression, Two Costumes.} Every radical can be written as an exponent, and every
rational exponent can be written as a radical --- and \emph{which costume you choose decides how
hard the problem is}. With your partner, translate carefully and say out loud which form you picked
and why. Unless a problem says otherwise, assume all variables represent \textbf{positive} real
numbers.
\end{headlinebox}

\vspace{0.06in}

\begin{tcolorbox}[colback=white, colframe=black!40, title=\textbf{Tier R --- Remediate},
    fonttitle=\bfseries, arc=1.5mm, left=3mm, right=3mm, top=2mm, bottom=2mm, breakable]
\textbf{1. Translate.} Fill in every empty cell. The \textbf{index} always becomes the
\textbf{denominator}.
\begin{center}
\renewcommand{\arraystretch}{1.9}
\begin{tabularx}{\linewidth}{@{}X X X@{}}
\hline
\textbf{Radical form} & \textbf{Rational exponent form} & \textbf{Value (numbers only)} \\
\hline
$\sqrt{x}$            & \ans{$x^{1/2}$} & --- \\
$\sqrt[3]{x^{2}}$     & \ans{$x^{2/3}$} & --- \\
\ans{$\sqrt{x^{5}}$}  & $x^{5/2}$     & --- \\
\ans{$\sqrt[6]{t}$}   & $t^{1/6}$     & --- \\
$\sqrt[4]{16}$        & \ans{$16^{1/4}$} & \ans{$2$} \\
\ans{$\sqrt[3]{8}$}   & $8^{1/3}$     & \ans{$2$} \\
$\left(\sqrt[3]{27}\right)^{2}$ & \ans{$27^{2/3}$} & \ans{$9$} \\
\hline
\end{tabularx}
\end{center}

\vspace{4pt}
\textbf{2. Where does each number go?} In the expression $a^{m/n}$:
\begin{itemize}[leftmargin=1.6em, itemsep=4pt]
  \item the number $n$ (the \emph{bottom}) is the \ans{index --- the root you take};
  \item the number $m$ (the \emph{top}) is the \ans{power you raise it to}.
\end{itemize}

\vspace{2pt}
\textbf{3. Evaluate.} Take the \emph{root} first, then the power --- the numbers stay small.
\begin{center}
\renewcommand{\arraystretch}{1.7}
\begin{tabular}{@{}l l l@{}}
(a) \; $9^{1/2}=$ \ans{$3$} & (b) \; $64^{1/3}=$ \ans{$4$} &
(c) \; $4^{3/2}=$ \ans{$8$} \\
(d) \; $25^{3/2}=$ \ans{$125$} & (e) \; $8^{4/3}=$ \ans{$16$} &
(f) \; $(-27)^{1/3}=$ \ans{$-3$} \\
\end{tabular}
\end{center}

\vspace{2pt}
\textbf{4.} One of the six above had a \emph{negative} base and still worked. Which one, and what is
it about its \textbf{denominator} that allowed it? \ansline{(f). The denominator is $3$, which is
\emph{odd}, so it asks for a cube root --- and cubing keeps the sign, so $(-3)^3=-27$. An even
denominator (a square or fourth root) would have failed on a negative base.}
\end{tcolorbox}

\begin{tcolorbox}[colback=white, colframe=black!40, title=\textbf{Tier A --- Approaching Proficiency},
    fonttitle=\bfseries, arc=1.5mm, left=3mm, right=3mm, top=2mm, bottom=2mm, breakable]
\textbf{1. No calculator.} Evaluate each. Show the root you took and the power you applied.
\begin{center}
\renewcommand{\arraystretch}{2.0}
\begin{tabularx}{\linewidth}{@{}X X X@{}}
(a) \; $27^{4/3}=$ \ans{$81$} &
(b) \; $16^{3/4}=$ \ans{$8$} &
(c) \; $(-8)^{2/3}=$ \ans{$4$} \\
(d) \; $81^{-1/2}=$ \ans{$\dfrac{1}{9}$} &
(e) \; $32^{-3/5}=$ \ans{$\dfrac{1}{8}$} &
(f) \; $\left(\dfrac{8}{27}\right)^{2/3}=$ \ans{$\dfrac{4}{9}$} \\
\end{tabularx}
\end{center}
\emph{Work:} (a) $\left(\sqrt[3]{27}\right)^4=3^4$. \; (b) $\left(\sqrt[4]{16}\right)^3=2^3$. \;
(c) $\left(\sqrt[3]{-8}\right)^2=(-2)^2$. \; (d) $\dfrac{1}{\sqrt{81}}$. \;
(e) $\dfrac{1}{\left(\sqrt[5]{32}\right)^3}=\dfrac{1}{2^3}$. \;
(f) $\left(\sqrt[3]{8/27}\right)^2=\left(\tfrac{2}{3}\right)^2$.

\vspace{3pt}
\textbf{2. Justify --- why root first?} Your partner insists on doing (b) as $\sqrt[4]{16^{3}}$.
That is not \emph{wrong}. Carry it out far enough to see what happens, then explain why taking the
root first is the better habit. \ansline{$16^3 = 4096$, so they now need $\sqrt[4]{4096}$ --- a
fourth root of a four-digit number, which almost nobody can do mentally. Root first keeps every
number small: $\sqrt[4]{16}=2$, then $2^3=8$. Both routes give $8$, because
$\sqrt[n]{a^m}=\left(\sqrt[n]{a}\right)^m$ --- but only one of them is arithmetic you can do in your
head.}

\vspace{2pt}
\textbf{3. Justify --- the sign trap.} A classmate says ``$81^{-1/2}$ must be negative, because the
exponent is negative.'' They are wrong. Explain what the negative exponent actually does, and give
the correct value. \ansline{A negative exponent means \emph{reciprocal}, not ``negative answer'':
$81^{-1/2}=\dfrac{1}{81^{1/2}}=\dfrac{1}{9}$. The minus sign moves the expression to the
denominator; it never changes the sign of the result. (A positive base raised to \emph{any} real
power stays positive.)}

\vspace{2pt}
\textbf{4. Sort them.} Without evaluating, decide whether each expression is a \emph{real number},
and give your reason (look at the denominator of the exponent and the sign of the base).
\begin{center}
\renewcommand{\arraystretch}{1.7}
\begin{tabularx}{\linewidth}{@{}l X X@{}}
\hline
\textbf{Expression} & \textbf{Real number? (yes/no)} & \textbf{Reason} \\
\hline
$(-64)^{1/3}$ & \ans{yes} & \ans{odd denominator: a cube root keeps the sign ($=-4$)} \\
$(-64)^{1/2}$ & \ans{no} & \ans{even denominator, negative base: nothing real squares to $-64$} \\
$(-32)^{1/5}$ & \ans{yes} & \ans{odd denominator: fifth root of a negative is negative ($=-2$)} \\
$(-16)^{3/4}$ & \ans{no} & \ans{even denominator, so the fourth root comes first --- and $\sqrt[4]{-16}$ is not real} \\
\hline
\end{tabularx}
\end{center}
\end{tcolorbox}

\begin{tcolorbox}[colback=white, colframe=black!40, title=\textbf{Tier E --- Extension},
    fonttitle=\bfseries, arc=1.5mm, left=3mm, right=3mm, top=2mm, bottom=2mm, breakable]
\textbf{Part 1 --- Problems radicals cannot do.} Each product or quotient below has \emph{mismatched
indices}, so no radical rule applies. Rewrite with rational exponents, simplify, and convert back to
a single radical.
\begin{center}
\renewcommand{\arraystretch}{2.1}
\begin{tabularx}{\linewidth}{@{}X X@{}}
(a) \; $\sqrt{x}\cdot\sqrt[4]{x} = $ \ans{$x^{3/4}=\sqrt[4]{x^{3}}$} &
(b) \; $\dfrac{\sqrt[3]{x^{2}}}{\sqrt{x}} = $ \ans{$x^{1/6}=\sqrt[6]{x}$} \\
(c) \; $\left(\sqrt[3]{x^{2}}\right)^{3/2} = $ \ans{$x^{1}=x$} &
(d) \; $\sqrt{\sqrt[3]{x}} = $ \ans{$x^{1/6}=\sqrt[6]{x}$} \\
\end{tabularx}
\end{center}
\emph{Work:} (a) $x^{1/2+1/4}=x^{3/4}$. \; (b) $x^{2/3-1/2}=x^{4/6-3/6}=x^{1/6}$. \;
(c) $x^{\frac{2}{3}\cdot\frac{3}{2}}=x^{1}$. \; (d) $\left(x^{1/3}\right)^{1/2}=x^{1/6}$.

\vspace{2pt}
For (d), explain what happened to the two indices and why that makes sense. \ansline{They
\emph{multiplied}: a cube root inside a square root is a sixth root, because
$\frac{1}{3}\cdot\frac{1}{2}=\frac{1}{6}$. That is just ``power to a power'' --- and it makes sense
because you are cutting $x$ into three equal factors and then cutting each of those in half.}

\vspace{5pt}
\textbf{Part 2 --- A rational exponent in the wild.} Kepler's third law says that a planet's
\textbf{orbital period} $T$ (in Earth years) depends on its \textbf{average distance} $a$ from the
Sun (in astronomical units, where $1$~AU is Earth's distance):
\[
  T(a) = a^{3/2}.
\]
\begin{center}
\renewcommand{\arraystretch}{1.25}
\begin{tabular}{|c|c|c|c|c|}
\hline
$a$ (AU)        & $1$ & $4$ & $9$ & $25$ \\ \hline
$T(a)$ (years)  & $1$ & \ans{$8$} & \ans{$27$} & \ans{$125$} \\ \hline
\end{tabular}
\end{center}
\begin{enumerate}[label=(\alph*), leftmargin=1.9em, itemsep=6pt]
  \item Complete the table \emph{without a calculator}. Which root did you take, and which power did
        you apply? \ansline{Square root first (denominator $2$), then cube it (numerator $3$):
        $4^{3/2}=\left(\sqrt{4}\right)^3=2^3=8$, and likewise $3^3=27$, $5^3=125$.}
  \item Write $T(a)=a^{3/2}$ in radical form: $T(a)=$ \ans{$\left(\sqrt{a}\right)^{3}$, or
        equivalently $\sqrt{a^{3}}$}. \; Which of the two
        equivalent radical forms did you use, and why is it the easier one to compute with?
        \ansline{$\left(\sqrt{a}\right)^3$ --- rooting first keeps the numbers small. The other form
        would ask for $\sqrt{15625}$ at $a=25$.}
  \item What is the domain of $T$ \emph{algebraically} (look at the denominator of the exponent)?
        \ans{$a\ge 0$, or $[0,\infty)$} \\
        Does the context of the problem restrict it any further? Explain. \ansline{Yes --- $a>0$.
        A distance of $0$~AU would put the planet inside the Sun, so the endpoint that algebra
        allows is not physically possible; the context throws it away.}
  \item Jupiter sits about $5.2$~AU from the Sun. Between which two \emph{table entries} must its
        orbital period fall, and roughly how many Earth years is that? \ansline{Between $a=4$ and
        $a=9$, so between $8$ and $27$ years --- and since $5.2$ is much closer to $4$, an estimate
        near $12$ years is reasonable. (Jupiter's actual period is about $11.9$ years.)}
  \item Solve $T=a^{3/2}$ for $a$ in terms of $T$. \quad $a = $ \ans{$T^{2/3}$} \\
        \emph{Hint:} raise both sides to the power that undoes $\frac{3}{2}$. Which power is that,
        and why does it work? \ansline{The reciprocal, $\frac{2}{3}$: $\left(a^{3/2}\right)^{2/3}
        =a^{\frac{3}{2}\cdot\frac{2}{3}}=a^{1}=a$. Multiplying by the reciprocal is what makes the
        exponent $1$.}
  \item Your answer to (e) undoes $T(a)$ --- it is the \textbf{inverse} relationship from Lesson 6.0,
        now written with exponents instead of graphs. Use it to check one entry from the table, and
        state what the exponent $\frac{2}{3}$ is doing in words. \ansline{Check: $27^{2/3}
        =\left(\sqrt[3]{27}\right)^2=3^2=9$~AU, which is the table's $a=9$ row. In words: take the
        \emph{cube root} of the period, then \emph{square} it --- distance from period, instead of
        period from distance.}
\end{enumerate}
\end{tcolorbox}

\begin{teachernote}
\textbf{Tier R} is pure translation and should be finished in under ten minutes; the only concept
question is item 4, which forces students to say ``odd \emph{denominator}'' rather than ``odd
index'' --- the vocabulary swap is the point.
\textbf{Tier A} item 2 is worth a whole-class share: let a group actually compute $16^3=4096$ on the
board before anyone rescues them. Item 3 kills the year's most durable exponent error. In item 4,
watch for students marking $(-16)^{3/4}$ real by cubing first ($(-16)^3$ is negative, and its fourth
root is still not real) --- the answer is the same either way, but the clean argument is that the
even denominator demands a fourth root of a negative number.
\textbf{Tier E} Part 1 is the lesson's thesis in four problems: none of (a)--(d) is doable with
radical rules alone. Part 2 is the inverse thread from Lesson 6.0 continued in exponent language,
and (e) is a genuine preview of Lesson 6.7 --- raising both sides to the reciprocal power is the
algebraic version of reflecting over $y=x$. Groups that finish early can be asked why the table's
first row reads $1\mapsto 1$: the units were \emph{defined} from Earth's own orbit, so $1^{3/2}=1$
is a statement about the choice of units, not about physics.
\end{teachernote}

\end{document}
